% math4cs-hw10-matching-flow-solution.tex

% !TEX program = xelatex
%%%%%%%%%%%%%%%%%%%%
% see http://mirrors.concertpass.com/tex-archive/macros/latex/contrib/tufte-latex/sample-handout.pdf
% for how to use tufte-handout
\documentclass[a4paper, justified]{tufte-handout}

\input{hw-preamble} % feel free to modify this file if you understand LaTeX well
%%%%%%%%%%%%%%%%%%%%
\title{计算机数学 (11-matching-flow: 匹配与网络流)}
\me{魏恒峰}{hfwei@hnu.edu.cn}{}{}
\date{2026年6月15日}
%%%%%%%%%%%%%%%%%%%%
\begin{document}
\maketitle
%%%%%%%%%%%%%%%%%%%%
\noplagiarism % PLEASE DON'T DELETE THIS LINE!
%%%%%%%%%%%%%%%%%%%%
\begin{abstract}
\end{abstract}
%%%%%%%%%%%%%%%%%%%%
\beginrequired
%%%%%%%%%%%%%%%

%%%%%%%%%%%%%%%
\begin{problem}
  设 $G = (X, Y, E)$ 是一个 $k$-正则 ($k > 0$) 二部图。
  请证明:
  \begin{enumerate}[(1)]
    \item $|X| = |Y|$;
    \item $G$ 有一个 $X$-完美匹配。
  \end{enumerate}
\end{problem}

\begin{proof}
  \begin{enumerate}[(1)]
    \item 因为 $G$ 是 $k$-正则二部图, 所以 $k |X| = k |Y|$, 故有 $|X| = |Y|$。
    \item 对于 $X$ 的任何子集 $S$, 从 $S$ 到 $N(S)$ 的边共有 $k|N|$ 条。
      这 $k|N|$ 条边均与 $N(S)$ 中的顶点相连, 所以 $k|N| \le k|N(S)|$,
      即有 $|N(S)| \ge |S|$。
      根据 Hall 定理, $G$ 存在 $X$-完美匹配。
  \end{enumerate}
\end{proof}
%%%%%%%%%%%%%%%

%%%%%%%%%%%%%%%
\begin{problem}
  请证明: 每个二部图 $G$ 都有一个大小 $\ge e(G)/\Delta(G)$ 的匹配~\footnote{$e(G)$表示$G$的边数。}。
  (提示: 使用 K\"{o}nig-Egerv\'{a}ry 定理。)
\end{problem}

\begin{proof}
  每个顶点最多覆盖 (cover) $\Delta(G)$ 条边。
  因此, $G$ 的最小点覆盖中包含 $\ge e(G)/\Delta(G)$ 个顶点。
  根据 K\"{o}nig-Egerv\'{a}ry 定理, $G$ 的最大匹配包含 $\ge e(G)/\Delta(G)$ 条边。
\end{proof}
%%%%%%%%%%%%%%%

%%%%%%%%%%%%%%%
\begin{problem}
  设 $Y$ 为集合,
  $\mathcal{A} = \set{A_{1}, \dots, A_{m}}$
  为包含 $m$ 个集合的集合, 其中 $A_{i} \subseteq Y$ (对 $1 \le i \le m$)。
  $\mathcal{A}$ 的相异代表系 (System of Distinct Representatives; SDR)
  是 $Y$ 中 $m$ 个不同元素 $a_{1}, \dots, a_{m}$ 构成的集合,
  其中 $a_{i} \in A_{i}$ (对 $1 \le i \le m$)。

  \noindent 请证明: $\mathcal{A}$ 有 SDR 当且仅当
  \[
    \forall S \subseteq \set{1, \dots, m}.\;
      \big\lvert \bigcup_{i \in S} A_{i} \big\rvert \ge |S|.
  \]
\end{problem}

\begin{proof}
  构造二部图 $G = (X, Y, E)$:
  其中, $X = \set{A_{1}, \dots, A_{m}}$,
  $Y = \set{a_{1}, \dots, a_{n}}$ (设 $|Y| = n$),
  并添加边 $\set{A_{i}, a_{j}} \in E$ 当且仅当 $a_{j} \in A_{i}$。
  因此, $\mathcal{A}$ 有 SDR 当且仅当 $G$ 有 $X$-完美匹配。

  已知
  \[
    \forall S \subseteq \set{1, \dots, m}.\;
      \big\lvert \bigcup_{i \in S} A_{i} \big\rvert \ge |S|,
  \]
  即 Hall 条件成立。因此, 根据 Hall 定理, $G$ 有 $X$-完美匹配。
  故, $\mathcal{A}$ 有 SDR。
\end{proof}
%%%%%%%%%%%%%%%

%%%%%%%%%%%%%%%
\begin{problem}
  请给出以下网络的一个最大流与一个最小割。
  要求给出 Ford-Fulkerson Method 运行过程。
  \fig{width = 0.50\textwidth}{figs/network-flow}
\end{problem}

\begin{proof}
  最大流如下图所示, 值为 $7$;
  最小割为 $(\set{s, a, b}, \set{c, d, e, t})$, 容量为 $7$。
  \fig{width = 0.50\textwidth}{figs/network-flow-value-7}

  Ford-Fulkerson Method 一种可能的运行过程如下:
  \begin{itemize}
    \item 初始化 $\forall e \in E.\; f(e) = 0$。
    \item 找到 $f$-增广路径: $s \to a \to d \to t$, 允许的最大增量为 2。
    \item 将 $f$ 增广为 $f_{1}$,
      其中 $f_{1}(s \to a) = f_{1}(a \to d) = f_{1}(d \to t) = 2$。
    \item 找到 $f_{1}$-增广路径: $s \to c \to e \to t$, 允许的最大增量为 4。
    \item 将 $f_{1}$ 增广为 $f_{2}$,
      其中 $f_{2}(s \to c) = f_{2}(c \to e) = f_{2}(e \to t) = 4$。
    \item 找到 $f_{2}$-增广路径: $s \to b \to d \to c \to e \to t$,
      允许的最大增量为 1。
    \item 将 $f_{2}$ 增广为 $f_{3}$,
      其中 $f_{3}(s \to b) = f_{3}(b \to d) = f_{3}(d \to c) = 1$,
      $f_{3}(c \to e) = f_{3}(e \to t) = 5$。
    \item 找不到 $f_{3}$-增广路径。因此, $f_{3}$ 为最大流。
      同时找到最小割 $(\set{s, a, b}, \set{c, d, e, t})$。
  \end{itemize}
\end{proof}
%%%%%%%%%%%%%%%

%%%%%%%%%%%%%%%
\begin{problem}
  考虑下面的定理:
  \begin{theorem}[不能告诉你名字的某个著名定理]
    设 $G = (V, E)$ 是无向连通图, $v, w \in V$ 是不同的两个顶点。
    则 $v$, $w$ 之间的边不相交的 (edge-disjoint)~\footnote{
      设 $P_{1}$, $P_{2}$ 是两条$v$、$w$间的路径。
      如果 $P_{1}$ 与 $P_{2}$ 没有公共边,
      则 $P_{1}$、$P_{2}$ 是 $v$, $w$ 之间的边不相交的路径。
    } 路径的最大条数
    等于最小$vw$-边割集~\footnote{
      设 $F \subseteq E$ 为集。
      如果 $G$ 删除 $F$ 后, $v$ 与 $w$ 不再连通,
      则称 $F$ 是 $vw$-边割集。
    }的大小。
  \end{theorem}

  \fig{width = 0.50\textwidth}{figs/edge-disjoint}

  \begin{enumerate}[(1)]
    \item 考虑图中的 $v$, $w$ 顶点。
      请给出 $v$、$w$ 间的一个最大边不相交的路径集合。
    \item 考虑图中的 $v$, $w$ 顶点。
      请给出一个最小的 $vw$-边割集。
    \item 请使用最大流-最小割定理证明上述定理~\footnote{恭喜!你刚刚证明了图论中的一个著名定理。}。
      \marginnote{这就是 Menger's theorem:
        \teal{\url{https://en.wikipedia.org/wiki/Menger\%27s\_theorem}}。
        如果直接证明该定理, 还是比较复杂的。
        使用最大流-最小割定理则容易很多。不过, 还是有些技术细节需要谨慎处理。}
  \end{enumerate}
\end{problem}

\begin{proof}
  \begin{enumerate}[(1)]
    \item 最大边不相交的路径集合为
      \[
        \set{v-p-s-u-w, v-q-s-x-w, v-q-t-y-w, v-r-t-z-w}.
      \]
      注意, $v-q$ 有两条重边。
    \item 最小的 $vw$-边割集为
      \[
        \set{ps, qs, qt, rt}.
      \]
    \item 首先将无向图 $G$ 转化成有向图 $G'$:
      将每条无向边 $uv$ 转化成两条有向边 $u \to v$ 与 $v \to u$。
      然后将每条有向边的容量设置为 1。
      \mfig{width = 0.80\textwidth}{figs/graph2digraph}

      考虑任意两个顶点 $v, w$。
      将最大边不相交的$vw$-路径的条数记为 $\lambda'(v, w)$,
      将最小 $vw$-边割集的大小记为 $\kappa'(v, w)$~\footnote{这是标准记法。}。
      由于任何一个 $vw$-边割集都要包含边不相交的每条$vw$-路径中的至少一条边, 所以
      \[
        \kappa'(v, w) \ge \lambda'(v, w).
      \]
      下面我们利用最大流-最小割定理证明 (记最大流为 $f$, 最小割为 $(S, T)$)
      \[
        \lambda'(v, w) \ge \max \text{val}(f)
                       = \min \text{cap}(S, T)
                       \ge \kappa'(v, w).
      \]

      经过上述转化, 可将图 $G'$ 视为源点为 $v$, 汇点为 $w$ 的网络。
      根据 Integrity Theorem~\footnote{我们在课上没有讲过这个定理。
      不过该定理很直观: 如果网络中每条边的容量都是整数,
      那么该网络存在每条边的流量都是整数的最大网络流。
      它的证明也很简单, 因为 Ford-Fulkerson Method 就可以保证这一点。},
      $G'$ 存在每条边的流量均为整数的最大网络流, 设为 $f$。
      如果存在某条无向边 $xy$, $f(x \to y) = f(y \to x) = 1$,
      则可以调整 $f$, 使得 $f(x \to y) = f(y \to x) = 0$。
      该调整不改变 $\text{val}(f)$。
      因此, 我们可以假设最大流 $f$ 中每条无向边 $xy$ 最多使用一次
      (即 $f(x \to y)$ 与 $f(y \to x)$ 不同时为 1)。
      所以, 该最大流 $f$ 包含了 $\text{val}(f)$ 条
      (两两)边不相交的 $vw$ 路径。这证明了
      \[
        \lambda'(v, w) \ge \max \text{val}(f).
      \]

      对于任意割 $(S, T)$,
      每条无向边 $xy$ 有且只有一个方向的有向边算作割的容量。因此,
      \[
        \text{cap}(S, T) = |[S, T]|.
      \]
      其中 $[S, T]$ 表示从 $S$ 到 $T$ 的有向边的集合。
      因为 $[S, T]$ 是一个边割集, 所以
      \[
        \min \text{cap}(S, T) \ge \kappa'(v, w).
      \]

      根据最大流-最小割定理,
      \[
        \lambda'(v, w) \ge \max \text{val}(f)
                       = \min \text{cap}(S, T)
                       \ge \kappa'(v, w).
      \]
  \end{enumerate}
\end{proof}
%%%%%%%%%%%%%%%
%%%%%%%%%%%%%%%%%%%%
\end{document}